\documentclass[lualatex,ja=standard]{bxjsarticle}
\usepackage{amsmath,amssymb}
\usepackage{enumitem}
\setlist[enumerate]{itemsep=0.6em, topsep=0.5em}
\pagestyle{empty}

\begin{document}

\noindent\textbf{数IA 思考演習（標準）}
\hfill\textbf{制限時間：25分}
\hrule
\vspace{0.8em}

\begin{enumerate}
\item $x$ を整数とする。
\[
|2x-5|\le 7
\]
を満たす $x$ をすべて求めよ。

\item $a$ を定数とする。2次方程式
\[
x^2-(a+1)x+a-2=0
\]
が異なる2つの実数解をもつような $a$ の範囲を求めよ。

\item $y=-x^2+4x+1$ について、
\begin{enumerate}[label=(\arabic*)]
\item 頂点の座標
\item $x$ 軸との共有点
\end{enumerate}
を求めよ。

\item 連立方程式
\[
\begin{cases}
x+y=5\\
x^2+y^2=13
\end{cases}
\]
を解け。

\item 三角形 $ABC$ で $AB=AC=10,\ BC=12$ である。
\begin{enumerate}[label=(\arabic*)]
\item $\angle A$ の大きさを求めよ。
\item 三角形の面積を求めよ。
\end{enumerate}

\item $0^\circ<\theta<180^\circ,\ \tan\theta=-\dfrac{4}{3}$ のとき、$\sin\theta,\cos\theta$ を求めよ。

\item 1から9までの数字が書かれたカード9枚から同時に2枚引く。
「和が偶数」である確率を求めよ。

\item 赤球4個、白球3個、黒球2個から同時に3個取り出す。
「3色すべて異なる色になる」取り出し方は何通りか。

\item あるクラス40人の数学テストの平均点は68点、中央値は70点だった。
次の記述のうち、必ず正しいものをすべて選べ。
\begin{itemize}
\item[(ア)] 70点以上の生徒は20人以上いる。
\item[(イ)] 68点以上の生徒は20人以上いる。
\item[(ウ)] 最高点は68点以上である。
\end{itemize}

\item 方程式
\[
\sqrt{x+5}=x-1
\]
を解け。
\end{enumerate}

\end{document}
